% STATUS: rewritten 2026-11-05
% Sources read in full: Kaygorodov 2007 (Theorems 2.5, 4.4 and the centroid
% remark in the introduction), Li 2018 (p. 657, Theorem 2.1, pp. 658--659).
\section{$\delta$-derivations and related maps}\label{sec:nonassociative:delta-quasi-anti}

Section~\ref{sec:nonassociative:alternative} introduced the
$\delta$-derivation maps and the alternative and Malcev cases
(Definition~\ref{def:nonassociative:delta-derivation}). Kaygorodov
treats the Jordan algebras and superalgebras
\cite{KacKaygorodov2007Delta}. For a simple finite-dimensional Jordan
algebra there is exactly one non-trivial value of $\delta$.

\begin{theorem}[{\cite[Theorem~2.5]{KacKaygorodov2007Delta}}]
\label{thm:nonassociative:delta-jordan}
Let $\varphi$ be a non-trivial $\delta$-derivation of a simple
finite-dimensional Jordan algebra $A$ over an algebraically closed field
$F$ of characteristic distinct from $2$. Then $\delta=\tfrac12$ and
$\varphi(x)=\alpha x$ for some $\alpha\in F$.
\end{theorem}

Kaygorodov notes that, for Lie algebras, a $\tfrac12$-derivation is an
element of the centroid of the algebra
\cite[Introduction]{KacKaygorodov2007Delta}. The Jordan superalgebra
case follows from the same classification method
\cite[Theorem~4.4]{KacKaygorodov2007Delta}.

\begin{theorem}[{\cite[Theorem~4.4]{KacKaygorodov2007Delta}}]
\label{thm:nonassociative:delta-jordan-super}
Let $A$ be a simple finite-dimensional Jordan superalgebra over an
algebraically closed field of characteristic $0$, and let $\varphi$ be a
non-trivial $\delta$-derivation of $A$. Then $\delta=\tfrac12$ and
$\varphi(x)=\alpha x$ for some $\alpha\in F$ and all $x\in A$.
\end{theorem}

The two operations of the next class exchange the roles of
multiplication and bracket in the Leibniz rule % claim-ok: own transition sentence, definitions cited below
\cite[p.~657]{AlbertiArestiLi2018Poisson}.

\begin{definition}[{\cite[p.~657]{AlbertiArestiLi2018Poisson}}]
\label{def:nonassociative:delta-poisson}
A \emph{Poisson algebra} is a vector space $B$ over a field $k$ with a
commutative associative multiplication $x\cdot y$ and a Lie bracket
$\{x,y\}$ satisfying
\[
\{x\cdot y,\,z\}\;=\;x\cdot\{y,z\}+\{x,z\}\cdot y,
\qquad x,y,z\in B.
\]
\end{definition}

The next definition states the derivations of a Poisson algebra.

\begin{definition}[{\cite[p.~657]{AlbertiArestiLi2018Poisson}}]
\label{def:nonassociative:delta-poisson-derivation}
A \emph{Poisson derivation} of a Poisson algebra $B$ is a linear map
$D$ on $B$ with $D(x\cdot y)=D(x)\cdot y+x\cdot D(y)$ and
$D(\{x,y\})=\{D(x),y\}+\{x,D(y)\}$ for all $x,y\in B$.
\end{definition}

The next example is the symplectic Poisson algebra.

\begin{example}[{\cite[p.~657]{AlbertiArestiLi2018Poisson}}]
\label{ex:nonassociative:symplectic-poisson}
The \emph{symplectic Poisson algebra} $S_n$ is the polynomial algebra
$k[x_1,y_1,\dots,x_n,y_n]$ with the bracket determined by
$\{x_i,y_j\}=\delta_{ij}$, $\{x_i,x_j\}=0$ and $\{y_i,y_j\}=0$.
\end{example}

On this algebra every derivation is inner.

\begin{theorem}[{\cite[Theorem~2.1,
pp.~658--659]{AlbertiArestiLi2018Poisson}}]
\label{thm:nonassociative:symplectic-inner}
Let $D$ be a derivation of the symplectic Poisson algebra
$S_n=k[x_1,y_1,\dots,x_n,y_n]$. Then $D^{*}=0$, and there exists
$f\in S_n$ such that $D=\operatorname{ad}_f$, that is,
$D(g)=\{f,g\}$ for every $g\in S_n$. Here $D^{*}$ denotes the
divergence of $D$ \cite[p.~657]{AlbertiArestiLi2018Poisson}.
\end{theorem}

Section~\ref{sec:nonassociative:poisson} develops the free side of the % claim-ok: own signpost to the next section
same picture.
